% GENERATED FILE. Edit canonical lessons, then run npm run generate:books.
% canonical-source-hash: fnv1a64:cca81936c55697df
% canonical-lessons: FA-C472-tabiat, FA-C472-hefazat-kardan, FA-C472-gardeshgar, FA-C472-atashfeshan, FA-C472-zistgah

\chapter{Protecting nature}
\label{ch:fa-protecting-nature}
\hlchaptermodality{472}

\begin{chapteropening}
\textbf{By the end of this chapter:} \emph{I can talk about looking after nature: nature, protecting it, tourists, volcanoes and habitats.}

Complete the last lesson of chapter 472: I can talk about looking after nature: nature, protecting it, tourists, volcanoes and habitats.
\end{chapteropening}

\section[\texorpdfstring{\fa{طبیعت} (tabi'at)}{طبیعت (tabi'at)}]{\fa{طبیعت} (tabi'at) — nature}

\label{lesson:FA-C472-tabiat}



\begin{quote}
Before the new one: say the Persian for a blossom, then the Persian for to bloom.
\end{quote}

\subsection*{You'll want to know: \fa{طبیعت}}
\textbf{\fa{طبیعت}} (\emph{tabi'at}) — \textquotedblleft{}nature\textquotedblright{}.

\begin{quote}
\fa{ما} \fa{جمعه} \fa{ها} \fa{به} \fa{دل} \fa{طبیعت} \fa{میرویم}.
\end{quote}

\begin{quote}
\emph{Mâ jom'e-hâ be del-e tabi'at miravim.}
\end{quote}

\begin{quote}
On Fridays we go out into the countryside.
\end{quote}

\emph{Del-e tabi'at} — the heart of nature, out in the wild.

\begin{cousinweb}
\emph{tabi'at} was borrowed from Arabic \emph{ṭabīʿa} 'nature, natural disposition' (EtymDB).
\end{cousinweb}

\begin{grammarlens}[title={What you've built}]
One more word for telling, explaining and describing at B1.
\end{grammarlens}

\subsection*{Your turn}
\begin{itemize}
\raggedright
  \item \emph{Say it:} \emph{tabi'at}
  \item \emph{Say it:} \emph{tabi'at}, once more
  \item \emph{Say it:} a sentence of your own with \emph{tabi'at}
  \item \emph{Recall:} say the Persian for a blossom, then the Persian for to bloom, then say \emph{tabi'at} again
\end{itemize}

\begin{tcolorbox}[breakable,colback=teal!4,colframe=teal!35!black,title={Before you move on}]
What does \fa{طبیعت} mean? (\textquotedblleft{}Nature\textquotedblright{}.) Say it once more.
\end{tcolorbox}
\section[\texorpdfstring{\fa{حفاظت} \fa{کردن} (hefâzat kardan)}{حفاظت کردن (hefâzat kardan)}]{\fa{حفاظت} \fa{کردن} (hefâzat kardan) — to protect, to conserve}

\label{lesson:FA-C472-hefazat-kardan}



\begin{quote}
Before the new one: say the Persian for to bloom, then the Persian for nature.
\end{quote}

\subsection*{You'll want to know: \fa{حفاظت} \fa{کردن}}
\textbf{\fa{حفاظت} \fa{کردن}} (\emph{hefâzat kardan}) — \textquotedblleft{}to protect, to conserve\textquotedblright{}.

\begin{quote}
\fa{مردم} \fa{روستا} \fa{از} \fa{جنگل} \fa{حفاظت} \fa{میکنند}.
\end{quote}

\begin{quote}
\emph{Mardom-e rustâ az jangal hefâzat mikonand.}
\end{quote}

\begin{quote}
The villagers protect the forest.
\end{quote}

The thing protected takes \emph{az}. \emph{Mantaqe-ye hefâzat-shode} — a protected area.

\begin{cousinweb}
\emph{hefâzat kardan} is built from \emph{hefâzat} 'protection', borrowed from Arabic \emph{ḥifāẓa}, which EtymDB derives from \emph{ḥafiẓa} 'to guard, to know by heart', the source of \emph{hefz kardan} 'to keep, to memorize', and \emph{kardan} 'to do'.
\end{cousinweb}

\begin{grammarlens}[title={What you've built}]
One more word for telling, explaining and describing at B1.
\end{grammarlens}

\subsection*{Your turn}
\begin{itemize}
\raggedright
  \item \emph{Say it:} \emph{hefâzat kardan}
  \item \emph{Say it:} \emph{hefâzat kardan}, once more
  \item \emph{Say it:} a sentence of your own with \emph{hefâzat kardan}
  \item \emph{Recall:} say the Persian for to bloom, then the Persian for nature, then say \emph{hefâzat kardan} again
\end{itemize}

\begin{tcolorbox}[breakable,colback=teal!4,colframe=teal!35!black,title={Before you move on}]
What does \fa{حفاظت} \fa{کردن} mean? (\textquotedblleft{}To protect, to conserve\textquotedblright{}.) Say it once more.
\end{tcolorbox}
\section[\texorpdfstring{\fa{گردشگر} (gardeshgar)}{گردشگر (gardeshgar)}]{\fa{گردشگر} (gardeshgar) — a tourist}

\label{lesson:FA-C472-gardeshgar}



\begin{quote}
Before the new one: say the Persian for nature, then the Persian for to protect.
\end{quote}

\subsection*{You'll want to know: \fa{گردشگر}}
\textbf{\fa{گردشگر}} (\emph{gardeshgar}) — \textquotedblleft{}a tourist\textquotedblright{}.

\begin{quote}
\fa{در} \fa{تابستان} \fa{این} \fa{جزیره} \fa{پر} \fa{از} \fa{گردشگر} \fa{است}.
\end{quote}

\begin{quote}
\emph{Dar tâbestân in jazire por az gardeshgar ast.}
\end{quote}

\begin{quote}
In summer this island is full of tourists.
\end{quote}

\emph{gardesh} 'outing' + \emph{-gar} 'doer'. Tourism is \emph{gardeshgari}.

\begin{cousinweb}
\emph{gardeshgar} is a native compound of \emph{gardesh} 'an outing', built on \emph{gashtan} 'to turn, to wander', and the ending \emph{-gar} 'doer'.
\end{cousinweb}

\begin{grammarlens}[title={What you've built}]
One more word for telling, explaining and describing at B1.
\end{grammarlens}

\subsection*{Your turn}
\begin{itemize}
\raggedright
  \item \emph{Say it:} \emph{gardeshgar}
  \item \emph{Say it:} \emph{gardeshgar}, once more
  \item \emph{Say it:} a sentence of your own with \emph{gardeshgar}
  \item \emph{Recall:} say the Persian for nature, then the Persian for to protect, then say \emph{gardeshgar} again
\end{itemize}

\begin{tcolorbox}[breakable,colback=teal!4,colframe=teal!35!black,title={Before you move on}]
What does \fa{گردشگر} mean? (\textquotedblleft{}A tourist\textquotedblright{}.) Say it once more.
\end{tcolorbox}
\section[\texorpdfstring{\fa{آتشفشان} (âtashfeshân)}{آتشفشان (âtashfeshân)}]{\fa{آتشفشان} (âtashfeshân) — a volcano}

\label{lesson:FA-C472-atashfeshan}



\begin{quote}
Before the new one: say the Persian for to protect, then the Persian for a tourist.
\end{quote}

\subsection*{You'll want to know: \fa{آتشفشان}}
\textbf{\fa{آتشفشان}} (\emph{âtashfeshân}) — \textquotedblleft{}a volcano\textquotedblright{}.

\begin{quote}
\fa{این} \fa{کوه} \fa{یک} \fa{آتشفشان} \fa{خاموش} \fa{است}.
\end{quote}

\begin{quote}
\emph{In kuh yek âtashfeshân-e khâmush ast.}
\end{quote}

\begin{quote}
This mountain is an extinct volcano.
\end{quote}

Literally 'fire-scatterer'.

\begin{cousinweb}
\emph{âtashfeshân} is a native compound of \emph{âtash} 'fire', inherited from Middle Persian \emph{ātaxš} (EtymDB), and \emph{feshân}, the stem of \emph{feshândan} 'to scatter'.
\end{cousinweb}

\begin{grammarlens}[title={What you've built}]
One more word for telling, explaining and describing at B1.
\end{grammarlens}

\subsection*{Your turn}
\begin{itemize}
\raggedright
  \item \emph{Say it:} \emph{âtashfeshân}
  \item \emph{Say it:} \emph{âtashfeshân}, once more
  \item \emph{Say it:} a sentence of your own with \emph{âtashfeshân}
  \item \emph{Recall:} say the Persian for to protect, then the Persian for a tourist, then say \emph{âtashfeshân} again
\end{itemize}

\begin{tcolorbox}[breakable,colback=teal!4,colframe=teal!35!black,title={Before you move on}]
What does \fa{آتشفشان} mean? (\textquotedblleft{}A volcano\textquotedblright{}.) Say it once more.
\end{tcolorbox}
\section[\texorpdfstring{\fa{زیستگاه} (zistgâh)}{زیستگاه (zistgâh)}]{\fa{زیستگاه} (zistgâh) — a habitat}

\label{lesson:FA-C472-zistgah}



\begin{quote}
Before the new one: say the Persian for a tourist, then the Persian for a volcano.
\end{quote}

\subsection*{You'll want to know: \fa{زیستگاه}}
\textbf{\fa{زیستگاه}} (\emph{zistgâh}) — \textquotedblleft{}a habitat\textquotedblright{}.

\begin{quote}
\fa{این} \fa{تالاب} \fa{زیستگاه} \fa{پرنده} \fa{های} \fa{زیادی} \fa{است}.
\end{quote}

\begin{quote}
\emph{In tâlâb zistgâh-e parande-hâ-ye ziyâdi ast.}
\end{quote}

\begin{quote}
This wetland is the habitat of many birds.
\end{quote}

\emph{zist} 'living' + \emph{gâh} 'place'.

\begin{cousinweb}
\emph{zistgâh} is a native compound of \emph{zist}, the past stem of \emph{zistan} 'to live', an inherited verb usually traced to Middle Persian \emph{zīstan}, and \emph{gâh} 'place'.
\end{cousinweb}

\begin{grammarlens}[title={What you've built}]
One more word for telling, explaining and describing at B1.
\end{grammarlens}

\subsection*{Your turn}
\begin{itemize}
\raggedright
  \item \emph{Say it:} \emph{zistgâh}
  \item \emph{Say it:} \emph{zistgâh}, once more
  \item \emph{Say it:} a sentence of your own with \emph{zistgâh}
  \item \emph{Recall:} say the Persian for a tourist, then the Persian for a volcano, then say \emph{zistgâh} again
\end{itemize}

\begin{tcolorbox}[breakable,colback=teal!4,colframe=teal!35!black,title={Before you move on}]
What does \fa{زیستگاه} mean? (\textquotedblleft{}A habitat\textquotedblright{}.) Say it once more.
\end{tcolorbox}
